\chapter{Adrenaline Autoinjector: How to Use It}

\section*{Definition and Overview}
An adrenaline autoinjector is a pen-shaped medical device that delivers a single, fixed dose of adrenaline into the outer thigh muscle to treat a severe allergic reaction (anaphylaxis). Common brands include EpiPen and EpiPen Jr, although the device design and instructions vary slightly between brands and countries, so users should always follow the specific instructions provided with their own device.

Autoinjectors are prescribed for people at risk of anaphylaxis, most often because of food allergy, insect-sting allergy, or allergy to medicines. The device is a first-aid measure: it rapidly reverses the most dangerous effects of the reaction and buys time, but it is not a substitute for emergency hospital care, and an ambulance must always be called whenever the device is used.

\section*{Epidemiology}
Anaphylaxis is a medical emergency that can affect people of any age. Around 1 to 2 in 100 people experience anaphylaxis at some point in life in some countries, and rates are rising in line with the increasing prevalence of food allergy, particularly in children. Common triggers include peanuts, tree nuts, shellfish, egg, cow's milk, insect stings (especially bees and wasps), and certain medicines.

Many people who carry an autoinjector have never had to use one, and studies show that when a reaction does occur, people often hesitate, delay, or make errors with the device. Practising regularly with a trainer device, keeping an in-date injector, and knowing the steps cold are the most effective ways to respond correctly under the stress of an emergency.

\section*{Aetiology and Pathophysiology}
Anaphylaxis is a rapid, severe, systemic allergic reaction. Within minutes of exposure to a trigger, the immune system floods the body with histamine and other chemical mediators released from mast cells and basophils.

\begin{itemize}
    \item \textbf{Airway:} swelling of the lips, tongue, and throat, which can narrow or obstruct breathing
    \item \textbf{Breathing:} bronchoconstriction causes wheezing and breathlessness
    \item \textbf{Circulation:} blood vessels widen and leak fluid into the tissues, causing a dangerous fall in blood pressure, dizziness, and collapse
    \item \textbf{Skin and gut:} hives and flushing, with vomiting, abdominal pain, or diarrhoea
    \item \textbf{How adrenaline works:} injected adrenaline narrows blood vessels (restoring blood pressure and reducing swelling), relaxes the airway muscles, and strengthens the heartbeat, usually within minutes of injection
\end{itemize}

\section*{Clinical Features}
Recognising the signs of anaphylaxis early is essential, because giving adrenaline promptly saves lives.

\begin{itemize}
    \item Hives, itching, flushing, or swollen skin
    \item Swelling of the lips, tongue, face, or throat
    \item Difficulty breathing, wheezing, a hoarse voice, or noisy breathing
    \item Vomiting, abdominal pain, or diarrhoea
    \item Dizziness, faintness, pale or clammy skin, and a rapid, weak pulse
    \item A sense of dread, or collapse and loss of consciousness
\end{itemize}

\section*{Differential Diagnosis}
In a suspected allergic emergency, it is safer to give adrenaline than to withhold it, because untreated anaphylaxis is life-threatening and the side effects of the drug are minor by comparison.

\begin{itemize}
    \item A vasovagal faint, which can look similar but does not cause breathing difficulty or swelling
    \item A severe asthma attack, which causes wheezing without skin or gut features
    \item A panic attack, which can cause rapid breathing and a racing heart
    \item Food poisoning or gastroenteritis, which can cause vomiting and collapse
    \item Choking, and low blood sugar (hypoglycaemia)
    \item In a person with a known severe allergy who has been exposed to their trigger, treat the episode as anaphylaxis without delay
\end{itemize}

\section*{When to Seek Medical Care}
Call for emergency help immediately if anaphylaxis is suspected, even if the autoinjector is used and symptoms improve. Adrenaline is first aid only, and hospital observation is always required because a second wave of symptoms can occur hours later.

\textbf{Call 000 (or local emergency services) for:}
\begin{itemize}
    \item Any breathing difficulty, throat swelling, or collapse after an allergic exposure
    \item Any use of an adrenaline autoinjector, even if the person feels better
    \item Worsening symptoms after the first dose has been given
    \item No improvement within five minutes of the first dose
\end{itemize}

\section*{Diagnosis and Investigations}
Anaphylaxis is diagnosed from the clinical features, and treatment must never be delayed by tests. Investigations belong to the period after the emergency.

\begin{itemize}
    \item \textbf{Clinical recognition:} the signs of anaphylaxis are the only "test" required before using the autoinjector
    \item \textbf{Hospital observation:} monitoring for a biphasic reaction, which can occur several hours after initial recovery
    \item \textbf{Allergy assessment:} referral to an allergy specialist for skin-prick testing or specific IgE blood tests to confirm the trigger after recovery
    \item \textbf{Action plan review:} updating the written anaphylaxis plan and device choice with the specialist
\end{itemize}

\section*{Management}
Correct management begins long before the emergency, with planning, storage, and regular practice.

\begin{itemize}
    \item \textbf{Plan ahead:} have a written anaphylaxis action plan and share it with family, school, and workplace contacts
    \item \textbf{Carry two autoinjectors:} many plans recommend carrying two, because a second dose may be needed
    \item \textbf{Store correctly:} keep the device at room temperature, away from heat, cold, and direct sunlight, and check the expiry date regularly
    \item \textbf{Replace before expiry:} an expired device may deliver a weaker dose or fail to work
    \item \textbf{Practise regularly:} use a trainer device so that using the real injector is automatic
    \item \textbf{After a reaction:} follow up with an allergy specialist to confirm the trigger and update the plan
\end{itemize}

When a reaction occurs, follow the specific instructions for your device and your written action plan. The general steps for a typical pen-style autoinjector are:

\begin{itemize}
    \item Hold the device firmly in your dominant hand, forming a fist around it
    \item Remove the safety cap (blue cap on the EpiPen) with the other hand; never place your thumb over the end of the device
    \item Place the tip of the device (orange end on the EpiPen) against the outer thigh; it can be injected through clothing
    \item Push firmly until you hear a click, and hold it in place for the time stated in the device instructions (commonly about 3 seconds for the EpiPen; check your own device)
    \item Remove the device and massage the injection site gently for about 10 seconds
    \item Note the time of the injection and record that it was given
    \item Call 000, keep the person lying flat with their legs raised (or sitting up if they are struggling to breathe), and watch them closely
    \item If there is no improvement within 5 minutes, give a second autoinjector if one is available
    \item Hand the used device to the ambulance crew, as it shows them how much adrenaline was given
\end{itemize}

\section*{Complications}
\begin{itemize}
    \item Not giving adrenaline early enough -- the greatest risk, and the reason to err on the side of using the device
    \item Accidental injection into a thumb or finger, which can cause intense pain and reduced blood flow and needs urgent medical review
    \item Short-lived adrenaline effects after a correct injection, including palpitations, tremor, and anxiety, which settle quickly
    \item A second wave of symptoms (biphasic reaction) hours after the first, which is why hospital observation is essential
    \item Failure to seek follow-up allergy assessment, leaving the trigger unidentified
\end{itemize}

\section*{Prognosis and Outcomes}
People who receive adrenaline promptly and are taken to hospital almost always recover fully from anaphylaxis. Timely use of an autoinjector greatly improves the chance of survival and a complete recovery. With careful trigger avoidance, a working in-date device, and a rehearsed plan, the risk of a serious outcome is very low.

\section*{Prevention and Self-Care}
\begin{itemize}
    \item \textbf{Carry everywhere:} carry your autoinjector at all times, and know exactly where it is
    \item \textbf{Avoid triggers:} avoid known allergens, and read food labels carefully
    \item \textbf{Inform others:} tell restaurants, schools, and workplaces about the allergy and how to use the device
    \item \textbf{Medical alert:} wear a medical alert bracelet describing the allergy
    \item \textbf{Check monthly:} inspect the expiry date and the appearance of the device each month
    \item \textbf{Practise:} practise with a trainer pen at least once a year
    \item \textbf{Review:} review your allergy action plan with your doctor regularly, especially before travelling
\end{itemize}

\section*{Sources and Further Reading}
The following sources provide reliable, up-to-date information on using adrenaline autoinjectors and managing anaphylaxis.

\begin{itemize}
    \item NHS. \textit{Conditions: anaphylaxis -- using an adrenaline autoinjector.} https://www.nhs.uk/conditions/
    \item Mayo Clinic. \textit{Anaphylaxis: first aid and when to use an epinephrine autoinjector.} https://www.mayoclinic.org
    \item MedlinePlus. \textit{Epinephrine injection: proper use and storage.} https://medlineplus.gov
    \item American College of Allergy, Asthma and Immunology. \textit{Anaphylaxis and autoinjector training.} https://acaai.org
    \item World Health Organization. \textit{Management of anaphylaxis in the community.} https://www.who.int
    \item MSD Manual. \textit{Anaphylaxis -- emergency treatment.} https://www.msdmanuals.com
\end{itemize}
